\section*{Chapter \ref{ch:b1:lewis-resonance-vsepr} --- Lewis Structures, Resonance and VSEPR}

\begin{solution}{exo:b1:lewis-resonance-vsepr:1}
\ce{H2O}: two \ce{O-H} bonds, two \omterm{def:b1:lewis-resonance-vsepr:lewis}{lone pairs} on O. \ce{NH3}: three
\ce{N-H} bonds, one \omterm{def:b1:lewis-resonance-vsepr:lewis}{lone pair} on N. \ce{CO2}: \ce{O=C=O}, two \omterm{def:b1:lewis-resonance-vsepr:lewis}{lone pairs}
on each O. \ce{HCN}: \ce{H-C#N}, one \omterm{def:b1:lewis-resonance-vsepr:lewis}{lone pair} on N. \ce{CH2O}: carbon
bonded to two H and double-bonded to O, two \omterm{def:b1:lewis-resonance-vsepr:lewis}{lone pairs} on O. All \omterm{def:b1:lewis-resonance-vsepr:formal-charge}{formal
charges} are zero in these five. \ce{NH4+}: four \ce{N-H} bonds, no \omterm{def:b1:lewis-resonance-vsepr:lewis}{lone
pair}; nitrogen $5 - 0 - 4 = +1$, the charge of the ion.
\end{solution}

\begin{solution}{exo:b1:lewis-resonance-vsepr:2}
\ce{H2S}: \ce{AX2E2}, bent, below $109.5^\circ$ (measured $92^\circ$).
\ce{PCl3}: \ce{AX3E}, trigonal pyramidal, about $100^\circ$.
\ce{BF3}: \ce{AX3}, trigonal planar, $120^\circ$. \ce{SiH4}: \ce{AX4},
tetrahedral, $109.5^\circ$. \ce{CO3^{2-}}: \ce{AX3} (the double bond
counts once), trigonal planar, $120^\circ$.
% ledger: geo:H2S.aHSH, geo:PCl3.aClPCl
\end{solution}

\begin{solution}{exo:b1:lewis-resonance-vsepr:3}
Non-polar: \ce{CCl4} (tetrahedral \ce{AX4}), \ce{BF3} (\ce{AX3}),
\ce{CO2} (linear \ce{AX2}): the \omterm{def:b1:lewis-resonance-vsepr:dipole}{bond dipoles} cancel. Polar: \ce{CH2Cl2}
(two different kinds of bonds on a tetrahedron), \ce{NH3} (pyramidal),
\ce{SO2} (bent).
\end{solution}

\begin{solution}{exo:b1:lewis-resonance-vsepr:4}
\ce{H+} is the \omterm{def:b1:lewis-resonance-vsepr:lewis-acid}{Lewis acid} and \ce{H2O} the base: the new \ce{O-H} bond of
\ce{H3O+} is dative when formed (afterwards the three \ce{O-H} bonds are
identical). \ce{Cu^{2+}} is the acid and each \ce{NH3} a base: the four
\ce{Cu-N} bonds are dative.
\end{solution}

\begin{solution}{exo:b1:lewis-resonance-vsepr:5}
\chemfig{CH_3-C(=[:60]O)-[:-60]\chemabove{O}{\scriptstyle\ominus}}
$\leftrightarrow$
\chemfig{CH_3-C(-[:60]\chemabove{O}{\scriptstyle\ominus})=[:-60]O}.
The two structures are equivalent: both \ce{C-O} bonds have order 1.5 and
the same length, between the double and the single bond. In methanoic
acid the two oxygens differ (one carries the H), the structures are not
equivalent, and the bonds keep distinct lengths, 120.2 (\ce{C=O}) and
\qty{134.3}{pm} (\ce{C-OH}).
\end{solution}

\begin{solution}{exo:b1:lewis-resonance-vsepr:6}
$N = 6 + 4 \times 6 + 2 = 32$ electrons. With four single bonds and three
\omterm{def:b1:lewis-resonance-vsepr:lewis}{lone pairs} on each oxygen: sulfur $6 - 0 - 4 = +2$, each oxygen
$6 - 6 - 1 = -1$; total $+2 - 4 = -2$. Shape \ce{AX4}, tetrahedral. With
two \ce{S=O} bonds, choosing which two of the four oxygens are
double-bonded gives $\binom{4}{2} = 6$ structures.
\end{solution}

\begin{solution}{exo:b1:lewis-resonance-vsepr:7}
\ce{SF4}: \ce{AX4E}, seesaw. \ce{ClF3}: \ce{AX3E2}, T-shaped.
\ce{XeF2}: \ce{AX2E3}, linear. \ce{XeF4}: \ce{AX4E2}, square planar. In a
trigonal bipyramid an axial position has three neighbours at $90^\circ$,
an equatorial position only two: the bulky \omterm{def:b1:lewis-resonance-vsepr:lewis}{lone pairs} go where they meet
the fewest $90^\circ$ repulsions.
\end{solution}

\begin{solution}{exo:b1:lewis-resonance-vsepr:8}
Sulfur and phosphorus are larger and less electronegative than oxygen and
nitrogen: the \omterm{def:b1:lewis-resonance-vsepr:lewis}{bonding pairs} lie farther from the central atom and are
pulled towards hydrogen, so they repel one another less, and the \omterm{def:b1:lewis-resonance-vsepr:lewis}{lone
pairs} squeeze the bonds towards $90^\circ$.
\end{solution}

\begin{solution}{exo:b1:lewis-resonance-vsepr:9}
\ce{CH3-C#N}: the \ce{CH3} carbon is sp$^3$, the nitrile carbon sp, the
nitrogen sp; 5 $\sigma$ bonds (three \ce{C-H}, \ce{C-C}, one in
\ce{C#N}) and 2 $\pi$ bonds. \ce{CH2=CH-CH3}: the two alkene carbons are
sp$^2$, the methyl carbon sp$^3$; 8 $\sigma$ bonds (six \ce{C-H}, two
\ce{C-C}) and 1 $\pi$ bond.
\end{solution}

\begin{solution}{exo:b1:lewis-resonance-vsepr:10}
Both molecules are pyramidal with a \omterm{def:b1:lewis-resonance-vsepr:lewis}{lone pair} on nitrogen. The \omterm{def:b1:lewis-resonance-vsepr:lewis}{lone pair}
gives a contribution pointing (negative to positive) from the \omterm{def:b1:lewis-resonance-vsepr:lewis}{lone pair}
towards the nucleus. In \ce{NH3}, nitrogen is the negative end of each
bond: the \omterm{def:b1:lewis-resonance-vsepr:dipole}{bond dipoles} point from N towards the hydrogens, in the same
direction as the lone-pair contribution, and they add. In \ce{NF3},
fluorine is the negative end: the \omterm{def:b1:lewis-resonance-vsepr:dipole}{bond dipoles} point from the fluorines
towards N, against the lone-pair contribution, and the two nearly cancel.
\end{solution}

\begin{solution}{exo:b1:lewis-resonance-vsepr:11}
$\mu_b = \mu/(2\cos(\theta/2)) = 1.63/(2\cos 59.75^\circ) =
\qty{1.62}{\debye} = \qty{5.40e-30}{C.m}$. Then $\delta = \mu_b/(e\,d) =
5.40 \times 10^{-30}/(1.602 \times 10^{-19} \times 143.2 \times 10^{-12})
= 0.24$.
\end{solution}

\begin{solution}{exo:b1:lewis-resonance-vsepr:12}
\ce{N=N=O} with \omterm{def:b1:lewis-resonance-vsepr:formal-charge}{formal charges} $-1$, $+1$, $0$; \ce{N#N-O} with $0$,
$+1$, $-1$. The \ce{N-N} bond (\qty{112.8}{pm}) is close to the triple
bond of \ce{N2} (\qty{109.8}{pm}): the structure \ce{N#N-O} weighs more,
and it places the negative charge on oxygen, the more electronegative
atom. The \ce{N-O} bond is however shorter than a single bond: the other
structure contributes too.
\end{solution}

\begin{solution}{pb:b1:lewis-resonance-vsepr:1}
\textbf{1.} \ce{N2} 10, \ce{NO} 11, \ce{NO2} 17, \ce{N2O} 16, \ce{O3} 18,
\ce{HNO3} 24.
\textbf{2.} \ce{N#N} with one \omterm{def:b1:lewis-resonance-vsepr:lewis}{lone pair} on each nitrogen.
\textbf{3.} \ce{N=O} with two \omterm{def:b1:lewis-resonance-vsepr:lewis}{lone pairs} on oxygen, one \omterm{def:b1:lewis-resonance-vsepr:lewis}{lone pair} and one
single electron on nitrogen. With 11 electrons, an odd number, one atom
always has an odd count: nitrogen has 7 electrons around it.
\textbf{4.} \ce{N=N=O}: $-1$, $+1$, $0$; \ce{N#N-O}: $0$, $+1$, $-1$.
\textbf{5.} \ce{O=O-O}: central oxygen $+1$ (one \omterm{def:b1:lewis-resonance-vsepr:lewis}{lone pair}, three bonds),
the singly bonded terminal oxygen $-1$, the other 0.
\textbf{6.} As in \cref{ex:b1:lewis-resonance-vsepr:hno3}: nitrogen $+1$,
the oxygen bearing a single bond and no hydrogen $-1$, the others 0.
\textbf{7.} \ce{NO} and \ce{NO2} (11 and 17 electrons).
\textbf{8.} Two equivalent structures, the double bond on either side:
order 1.5 for each bond.
\textbf{9.} Yes: \qty{127.8}{pm} lies between the double bond
(\qty{120.8}{pm}) and the single bond (\qty{147.5}{pm}).
\textbf{10.} The double bond on each of the three oxygens in turn; each
\ce{N-O} bond is double in one structure of three: order $4/3$.
\textbf{11.} The long bond, \qty{140.6}{pm}, is \ce{N-OH}, a single bond.
The two oxygens without hydrogen share one $\pi$ bond through two
equivalent \omterm{def:b1:lewis-resonance-vsepr:resonance}{resonance structures}: two bonds of order 1.5, 121.1 and
\qty{119.9}{pm}.
\textbf{12.} Nitrogen carries a single electron rather than a \omterm{def:b1:lewis-resonance-vsepr:lewis}{lone pair}:
it repels the \omterm{def:b1:lewis-resonance-vsepr:lewis}{bonding pairs} less than a pair would, and the bonds open
beyond $120^\circ$.
\textbf{13.} \ce{N#N-O}, with the negative charge on oxygen.
\textbf{14.} \ce{O3}: \ce{AX2E}, bent. \ce{N2O}: \ce{AX2} (central N),
linear. \ce{NO3-}: \ce{AX3}, trigonal planar. Nitrogen of \ce{HNO3}:
\ce{AX3}, trigonal planar.
\textbf{15.} The \omterm{def:b1:lewis-resonance-vsepr:lewis}{lone pair} of the central oxygen repels the \omterm{def:b1:lewis-resonance-vsepr:lewis}{bonding pairs}
more than they repel each other.
\textbf{16.} Ozone is bent and its central atom differs from the end
atoms (positive \omterm{def:b1:lewis-resonance-vsepr:formal-charge}{formal charge} in the middle, negative charge shared by
the ends): a small dipole along the bisector. \ce{N2O} is linear but its
two ends are different atoms: a small dipole. \ce{N2} is symmetric: none.
\textbf{17.} \ce{CO2} is linear (\ce{AX2}): its \omterm{def:b1:lewis-resonance-vsepr:dipole}{bond dipoles} cancel.
\ce{SO2} is bent (\ce{AX2E}): they add.
\textbf{18.} \ce{AX3E}: the \omterm{def:b1:lewis-resonance-vsepr:lewis}{lone pair} closes the angles below
$109.5^\circ$; measured $106.7^\circ$.
\textbf{19.} \ce{AX2E2}: two \omterm{def:b1:lewis-resonance-vsepr:lewis}{lone pairs} close it further; $104.5^\circ$.
\textbf{20.} $\mu = 2\mu_{OH}\cos(\theta/2)$
(\cref{prop:b1:lewis-resonance-vsepr:dipole-sum}).
\textbf{21.} $\mu_{OH} = 1.857/(2\cos 52.24^\circ) = 1.857/1.2246 =
\qty{1.52}{\debye}$.
\textbf{22.} $1.516 \times 3.336 \times 10^{-30} = \qty{5.06e-30}{C.m}$.
\textbf{23.} $e\,d = 1.602 \times 10^{-19} \times 95.8 \times 10^{-12} =
\qty{1.535e-29}{C.m}$, that is \qty{4.60}{\debye}.
\textbf{24.} $\delta = 5.06 \times 10^{-30}/1.535 \times 10^{-29} =
\textbf{0.33}$: the \ce{O-H} bond of water carries a third of the charge
a fully ionic bond would carry.
\end{solution}
